\section{Methodology}
\label{sec:method}

\subsection{Problem Formulation}
Let $x \in \mathbb{R}^8$ denote the measurement vector -- ammonia, BOD, dissolved oxygen,
orthophosphate, pH, temperature, nitrogen and nitrate -- and $y \in [0,100]$ the CCME
index. The index is defined~\cite{ccme2001wqi} from a scope factor $F_1$, a frequency
factor $F_2$ and an amplitude factor $F_3$ as
\begin{equation}
    y = 100 - \frac{\sqrt{F_1^2 + F_2^2 + F_3^2}}{1.732},
    \label{eq:ccme}
\end{equation}
where $\sqrt{3} \approx 1.732$ normalises the maximum attainable distance to 100. Because
each record in this corpus is an individual measurement rather than an aggregated time
window, the upstream pipeline sets $F_2 \equiv F_1$; we note this adaptation because it
departs from the protocol as published and affects how the index should be read.

Equation~\eqref{eq:ccme} is deterministic in $x$. Learning $f : x \mapsto y$ from the
complete vector is therefore function approximation, and its error is bounded below only
by the capacity of the approximator. We instead study
\begin{equation}
    f_S : x_S \mapsto y, \qquad S \subset \{1,\dots,8\},
    \label{eq:partial}
\end{equation}
where $x_S$ is the measurement vector restricted to an available subset. For
$|S| < 8$ the formula cannot be evaluated, no closed-form solution exists, and the model
must exploit statistical structure among parameters, sites and water body types. We
report $|S| \in \{8,7,6,4\}$ and retain $|S| = 8$ as an emulation benchmark rather than a
prediction result.

The subsets are chosen to be interpretable rather than exhaustive. The 7/8 regime
withholds orthophosphate, which Section~\ref{sec:results} shows to be the dominant driver
of the index; 6/8 additionally withholds BOD, removing the heavy-organics signal; and 4/8
retains only the core physical-chemical suite of pH, dissolved oxygen, temperature and
nitrate, approximating what a continuously deployed in-situ sensor package measures
without laboratory assay.

\subsection{Splitting}
\label{sec:split}
All experiments share one partition, generated once and persisted to disk. Every model at
every availability level is validated and tested on byte-identical records, so reported
differences between models are not attributable to differences in evaluation data. The
training sets differ in size, and we state that explicitly in Section~\ref{sec:models}
rather than claiming a parity we do not have.

Records are grouped by monitoring station (\texttt{Area}, 61{,}689 distinct sites) and
whole groups are allocated to training, validation or test in a $70/15/15$ ratio, with
allocation stratified by country. No site appears in more than one partition. This
follows directly from the concern raised in~\cite{roberts2017cv}: the same station
sampled on consecutive dates yields near-identical chemistry, so a record-level random
split would place near-duplicates on both sides and report a figure that reflects site
memorisation rather than generalisation. Within each country, sites are shuffled under a
fixed seed and assigned by greedy bin-packing in descending size order, each site going
to whichever partition has the largest remaining deficit, which matches the target
proportions as closely as integer site sizes permit.

One country resists this treatment. China contributes a single site (Hou Bay), which
cannot be divided without breaking the zero-leakage guarantee, and is therefore allocated
entirely to training. Chinese records consequently appear in the reported per-country
breakdown as in-sample and are labelled as such; they are not evidence of generalisation.

\subsection{Preprocessing and Features}
The target is isolated before any transformation and never enters the design matrix. All
fitted statistics -- imputer medians, category encodings and scalers -- are learned from
the training partition alone and applied unchanged to validation and test. Masking for a
partial-measurement regime removes the withheld columns and, importantly, any derived
feature computed from them, so that no masked measurement re-enters through a ratio or a
transform.

The raw feature set comprises the available measurements together with country and water
body type. Site identifiers are excluded: with 61{,}689 levels under a site-grouped
split, every test site is unseen by construction, so the variable can only act as noise.
The engineered set adds logarithmic transforms of the skewed measurements -- ammonia,
BOD, nitrate, nitrogen and orthophosphate, whose skewness reaches 11.98 -- and ratios
between chemical parameters. Flags marking upstream-imputed values
(Section~\ref{sec:imputation}) are withheld from the primary feature set and evaluated
separately by ablation, because they are confounded with country.

Programmatic checks assert, for every configuration, that the target does not appear
among the features and that no group spans two partitions.

\subsection{Models and Training Budgets}
\label{sec:models}
Three predictors are compared. A median predictor establishes the floor any model must
clear. Random Forest is tuned over tree count, depth, minimum leaf size and features per
split. XGBoost is tuned over learning rate, depth, subsample, column subsample, minimum
child weight and $L_1$/$L_2$ regularisation, with early stopping on the validation
partition.

The two ensembles are not trained on the same number of records. The partition assigns
1{,}986{,}397 records to training, and XGBoost uses all of them; Random Forest is fitted
on a uniformly drawn subsample of 100{,}000, approximately 5\% of the available training
data. The reason is computational: Random Forest grows unpruned trees whose memory and
fitting cost scale with the number of records retained at every node, and fitting the
full partition was not feasible within our budget, whereas boosting's sequential,
depth-limited trees are far cheaper per record.

This is a genuine asymmetry and we do not want it read as a like-for-like comparison. The
two models see identical validation and test records, identical preprocessing and the
same feature sets, but unequal training data. The reported margin in favour of XGBoost
therefore reflects a combination of algorithmic difference and training budget, and part
of the gap would likely close if Random Forest were given the full partition. We return
to this in Section~\ref{sec:limitations}. Both expose the same training interface so that the Phase 2 neural model can
enter the identical evaluation path. Hyperparameters are selected on validation only; the
test partition is evaluated once, at the end.

\subsection{Metrics}
We report mean absolute error, root mean squared error and the coefficient of
determination, computed through a single shared implementation so that every model is
scored identically. MAE is reported alongside RMSE because the two respond differently to
the heavy tails documented in Section~\ref{sec:data}: RMSE is dominated by the small
number of large excursions, while MAE describes the typical record. $R^2$ is reported for
comparability with the water quality literature, with the caveat that its baseline is the
target variance, which differs across availability regimes.

Confidence intervals are obtained by bootstrap resampling of the test partition, which
makes the comparison between models a statement about sampling variability rather than a
single point estimate. As noted in Section~\ref{sec:ceiling}, the target is bounded above
and carries a large point mass at that bound; the metrics reported here are pooled across
it, and Section~\ref{sec:limitations} states what that conflates.
